\documentclass[12pt, twoside]{article}
\input{../preamble}

\fancyhead[LO]{La Scuola d'Italia / Dr.\ Huson / IB Math \\* 30 March 2026}
\fancyhead[RO]{Markscheme}
\fancyfoot[C]{\thepage}

\begin{document}
\subsubsection*{6.3 Classwork: Linear and Quadratic Functions --- Markscheme}

\begin{enumerate}[itemsep=0.2cm]

%--- Problem 1: Line through two points [8 marks] ---
\item A straight line $L$ passes through the points $A(-2, 5)$ and $B(4, -1)$.

\begin{enumerate}[itemsep=0.3cm]
  \item Find the gradient of $L$. \hfill [2]

  \begin{minipage}[t]{0.62\textwidth}
  $m = \dfrac{-1 - 5}{4 - (-2)} = \dfrac{-6}{6}$\\[0.15cm]
  $= -1$
  \end{minipage}%
  \begin{minipage}[t]{0.35\textwidth}\raggedleft
  \textbf{(M1)}\\[0.15cm]
  \textbf{(A1)}
  \end{minipage}

  \medskip
  \item Find the equation of $L$ in the form $y = mx + c$. \hfill [2]

  \begin{minipage}[t]{0.62\textwidth}
  $y - 5 = -1(x - (-2))$\\[0.15cm]
  $y = -x + 3$
  \end{minipage}%
  \begin{minipage}[t]{0.35\textwidth}\raggedleft
  \textbf{(M1)}\\[0.15cm]
  \textbf{(A1)}
  \end{minipage}

  \medskip
  \item Find the $x$-intercept of $L$. \hfill [2]

  \begin{minipage}[t]{0.62\textwidth}
  $0 = -x + 3$\\[0.15cm]
  $x = 3$ \quad $\Rightarrow (3,\; 0)$
  \end{minipage}%
  \begin{minipage}[t]{0.35\textwidth}\raggedleft
  \textbf{(M1)}\\[0.15cm]
  \textbf{(A1)}
  \end{minipage}

  \medskip
  \item A second line $L_2$ is parallel to $L$ and passes through the
        point $(0, 7)$. Find the equation of $L_2$. \hfill [2]

  \begin{minipage}[t]{0.62\textwidth}
  Parallel $\Rightarrow m = -1$\\[0.15cm]
  $y = -x + 7$
  \end{minipage}%
  \begin{minipage}[t]{0.35\textwidth}\raggedleft
  \textbf{(M1)}\\[0.15cm]
  \textbf{(A1)}
  \end{minipage}
\end{enumerate}

\newpage

%--- Problem 2: Factoring and properties of quadratic [8 marks] ---
\item Let $f(x) = x^2 - 4x - 12$.

\begin{enumerate}[itemsep=0.3cm]
  \item Factorise $f(x)$. \hfill [2]

  \begin{minipage}[t]{0.62\textwidth}
  $f(x) = (x - 6)(x + 2)$
  \end{minipage}%
  \begin{minipage}[t]{0.35\textwidth}\raggedleft
  \textbf{(A1)(A1)}
  \end{minipage}

  \medskip
  \item Hence find the $x$-intercepts of the graph of $f$. \hfill [1]

  \begin{minipage}[t]{0.62\textwidth}
  $x = 6$ and $x = -2$
  \end{minipage}%
  \begin{minipage}[t]{0.35\textwidth}\raggedleft
  \textbf{(A1)}
  \end{minipage}

  \medskip
  \item Find the coordinates of the vertex of the parabola. \hfill [3]

  \begin{minipage}[t]{0.62\textwidth}
  $x = \dfrac{-(-4)}{2(1)} = 2$\\[0.15cm]
  $f(2) = 4 - 8 - 12 = -16$\\[0.15cm]
  Vertex $= (2,\; -16)$
  \end{minipage}%
  \begin{minipage}[t]{0.35\textwidth}\raggedleft
  \textbf{(M1)}\\[0.15cm]
  \textbf{(A1)}\\[0.15cm]
  \textbf{(A1)}
  \end{minipage}

  \medskip
  \item Write down the $y$-intercept. \hfill [1]

  \begin{minipage}[t]{0.62\textwidth}
  $f(0) = -12$ \quad $\Rightarrow (0,\; -12)$
  \end{minipage}%
  \begin{minipage}[t]{0.35\textwidth}\raggedleft
  \textbf{(A1)}
  \end{minipage}

  \medskip
  \item Write down the range of $f$. \hfill [1]

  \begin{minipage}[t]{0.62\textwidth}
  $f(x) \geq -16$
  \end{minipage}%
  \begin{minipage}[t]{0.35\textwidth}\raggedleft
  \textbf{(A1)}
  \end{minipage}
\end{enumerate}

\newpage

%--- Problem 3: Line forms and perpendicular [9 marks] ---
\item The line $L_1$ has equation $2x + y - 8 = 0$.

\begin{enumerate}[itemsep=0.3cm]
  \item Write the equation of $L_1$ in the form $y = mx + c$. \hfill [2]

  \begin{minipage}[t]{0.62\textwidth}
  $2x + y - 8 = 0$\\[0.15cm]
  $y = -2x + 8$
  \end{minipage}%
  \begin{minipage}[t]{0.35\textwidth}\raggedleft
  \textbf{(M1)}\\[0.15cm]
  \textbf{(A1)}
  \end{minipage}

  \medskip
  \item Write down the gradient and the $y$-intercept of $L_1$. \hfill [1]

  \begin{minipage}[t]{0.62\textwidth}
  Gradient $= -2$, \quad $y$-intercept $= 8$
  \end{minipage}%
  \begin{minipage}[t]{0.35\textwidth}\raggedleft
  \textbf{(A1)}
  \end{minipage}

  \medskip
  \item The line $L_2$ is perpendicular to $L_1$ and passes through the
        point $(2, -1)$. Find the equation of $L_2$ in the form
        $ax + by + d = 0$, where $a$, $b$, $d \in \mathbb{Z}$. \hfill [3]

  \begin{minipage}[t]{0.62\textwidth}
  Perpendicular gradient: $m_2 = \dfrac{1}{2}$\\[0.2cm]
  $y - (-1) = \dfrac{1}{2}(x - 2)$\\[0.15cm]
  $2y + 2 = x - 2$\\[0.15cm]
  $x - 2y - 4 = 0$
  \end{minipage}%
  \begin{minipage}[t]{0.35\textwidth}\raggedleft
  \textbf{(A1)}\\[0.2cm]
  \textbf{(M1)}\\[0.15cm]
  ~\\[0.15cm]
  \textbf{(A1)}
  \end{minipage}

  \medskip
  \item Find the coordinates of the point of intersection of $L_1$
        and $L_2$. \hfill [3]

  \begin{minipage}[t]{0.62\textwidth}
  $-2x + 8 = \dfrac{1}{2}x - 2$\\[0.15cm]
  $-4x + 16 = x - 4$\\[0.15cm]
  $5x = 20 \Rightarrow x = 4$, \quad $y = -2(4) + 8 = 0$\\[0.15cm]
  Point of intersection: $(4,\; 0)$
  \end{minipage}%
  \begin{minipage}[t]{0.35\textwidth}\raggedleft
  \textbf{(M1)}\\[0.15cm]
  \textbf{(A1)}\\[0.15cm]
  ~\\[0.15cm]
  \textbf{(A1)}
  \end{minipage}
\end{enumerate}

\newpage

%--- Problem 4: Vertex form and reflection [8 marks] ---
\item Let $g(x) = 2x^2 - 8x + 6$.

\begin{enumerate}[itemsep=0.3cm]
  \item Write $g(x)$ in the form $a(x - h)^2 + k$. \hfill [3]

  \begin{minipage}[t]{0.62\textwidth}
  $g(x) = 2(x^2 - 4x) + 6$\\[0.15cm]
  $= 2(x^2 - 4x + 4) - 8 + 6$\\[0.15cm]
  $= 2(x - 2)^2 - 2$
  \end{minipage}%
  \begin{minipage}[t]{0.35\textwidth}\raggedleft
  \textbf{(M1)}\\[0.15cm]
  \textbf{(A1)}\\[0.15cm]
  \textbf{(A1)}
  \end{minipage}

  \medskip
  \item Write down the coordinates of the vertex of the parabola. \hfill [1]

  \begin{minipage}[t]{0.62\textwidth}
  Vertex $= (2,\; -2)$
  \end{minipage}%
  \begin{minipage}[t]{0.35\textwidth}\raggedleft
  \textbf{(A1)}
  \end{minipage}

  \medskip
  \item Find the $x$-intercepts of the graph of $g$. \hfill [2]

  \begin{minipage}[t]{0.62\textwidth}
  $2(x - 2)^2 - 2 = 0$\\[0.15cm]
  $(x - 2)^2 = 1 \Rightarrow x - 2 = \pm 1$\\[0.15cm]
  $x = 1$ and $x = 3$
  \end{minipage}%
  \begin{minipage}[t]{0.35\textwidth}\raggedleft
  \textbf{(M1)}\\[0.15cm]
  ~\\[0.15cm]
  \textbf{(A1)}
  \end{minipage}

  \medskip
  \item The graph of $g$ is reflected in the $x$-axis. Write the equation
        of the new graph in the form $y = ax^2 + bx + c$. \hfill [2]

  \begin{minipage}[t]{0.62\textwidth}
  $y = -g(x) = -(2x^2 - 8x + 6)$\\[0.15cm]
  $y = -2x^2 + 8x - 6$
  \end{minipage}%
  \begin{minipage}[t]{0.35\textwidth}\raggedleft
  \textbf{(M1)}\\[0.15cm]
  \textbf{(A1)}
  \end{minipage}
\end{enumerate}

\end{enumerate}
\end{document}
